% =====================================================================
\begin{thebibliography}{99}
\addcontentsline{toc}{chapter}{Bibliografía}
\markboth{Bibliografía}{}
\small

\bibitem{apostol}
T.~M. Apostol.
\newblock \emph{Calculus}, vol.~1: \emph{One-Variable Calculus, with an Introduction to Linear Algebra}, 2.ª ed.
\newblock John Wiley \& Sons, 1967.
\newblock Caps.~1 «The Concepts of Integral Calculus», 5 «The Relation between Integration and Differentiation», 7 «Polynomial Approximations to Functions», 10 «Sequences, Infinite Series, and Improper Integrals» y 11 «Sequences and Series of Functions».

\bibitem{mit1801}
D.~Jerison \emph{et al.}
\newblock 18.01SC: Single Variable Calculus.
\newblock Massachusetts Institute of Technology, MIT OpenCourseWare, otoño de 2010. Unidad~3 «The Definite Integral and its Applications», unidad~4 «Techniques of Integration» y unidad~5 «Exploring the Infinite»; colecciones de ejercicios \emph{18.01 Exercises}, unidad~3 «Integration» (secciones 3A--3G) y unidad~5 «Integration Techniques» (secciones 5A--5F).
\newblock \url{https://ocw.mit.edu/courses/18-01sc-single-variable-calculus-fall-2010/}

\bibitem{openstax1}
G.~Strang y E.~Herman.
\newblock \emph{Calculus Volume~1}.
\newblock OpenStax, Rice University, 2016. Cap.~5 «Integration» (§5.1--5.7).
\newblock \url{https://openstax.org/details/books/calculus-volume-1}

\bibitem{openstax2}
G.~Strang y E.~Herman.
\newblock \emph{Calculus Volume~2}.
\newblock OpenStax, Rice University, 2016. Caps.~3 «Techniques of Integration» (§3.1--3.7), 5 «Sequences and Series» (§5.1--5.6) y 6 «Power Series» (§6.1--6.4).
\newblock \url{https://openstax.org/details/books/calculus-volume-2}

\bibitem{spivak}
M.~Spivak.
\newblock \emph{Calculus}, 4.ª ed.
\newblock Publish or Perish, 2008.
\newblock Caps.~13 «Integrals», 14 «The Fundamental Theorem of Calculus», 18 «The Logarithm and Exponential Functions», 19 «Integration in Elementary Terms», 20 «Approximation by Polynomial Functions», 22 «Infinite Sequences», 23 «Infinite Series» y 24 «Uniform Convergence and Power Series».

\bibitem{stewart}
J.~Stewart.
\newblock \emph{Calculus: Early Transcendentals}, 8.ª ed.
\newblock Cengage Learning, 2015.

\bibitem{strang}
G.~Strang.
\newblock \emph{Calculus}, 3.ª ed.
\newblock Wellesley-Cambridge Press, 2017. Edición en línea en MIT OpenCourseWare, RES.18-001, otoño de 2023. Caps.~5 «Integrals», 7 «Techniques of Integration», 10 «Infinite Series» y §14.2 «Changing to Better Coordinates».
\newblock \url{https://ocw.mit.edu/courses/res-18-001-calculus-fall-2023/}

\end{thebibliography}
